\section{Robustness Tables}
\label{app:online}

This appendix collects auxiliary tables referenced in the main body.
All quantities are computed from the same pipeline and seed
documented in Section~\ref{app:reproducibility}.

\subsection{Cluster-gap-threshold robustness for H1}
\label{app:cluster-gap}

Section~\ref{ssec:clustering} treats the 15-minute cluster-gap
threshold as a design parameter and notes that the central H1 result
is insensitive to the choice within the range 5--30 minutes.
Table~\ref{tab:cluster-gap-robustness} reports the H1 event/baseline
magnitude ratio for the four thresholds. The ratio is monotone in
the threshold (tighter thresholds retain more independent clusters
and yield slightly higher ratios), but remains above $1.0$ in every
cell, the central conclusion in Section~\ref{ssec:results-volatility}.

\begin{table}[!htbp]
\centering
\caption{Cluster-gap-threshold robustness for the H1 magnitude ratio.}
\label{tab:cluster-gap-robustness}
\begin{threeparttable}
\small
\begin{tabular}{lccccc}
\toprule
 & & \multicolumn{4}{c}{Cluster gap threshold} \\
\cmidrule(lr){3-6}
Asset & Window & 5 min & 10 min & 15 min & 30 min \\
\midrule
EUR/USD & +5m & 1.90$\times$ & 1.85$\times$ & 1.76$\times$ & 1.67$\times$ \\
EUR/USD & +15m & 1.83$\times$ & 1.77$\times$ & 1.69$\times$ & 1.60$\times$ \\
EUR/USD & +60m & 1.66$\times$ & 1.62$\times$ & 1.58$\times$ & 1.50$\times$ \\
NDX & +5m & 1.73$\times$ & 1.65$\times$ & 1.56$\times$ & 1.48$\times$ \\
NDX & +15m & 1.71$\times$ & 1.66$\times$ & 1.57$\times$ & 1.47$\times$ \\
NDX & +60m & 1.57$\times$ & 1.51$\times$ & 1.43$\times$ & 1.38$\times$ \\
\addlinespace
\multicolumn{6}{l}{\textit{Number of clusters after dedupe}} \\
\quad EUR/USD +5m & & 1,557 & 1,395 & 1,288 & 1,069 \\
\quad EUR/USD +15m & & 1,539 & 1,379 & 1,273 & 1,057 \\
\quad EUR/USD +60m & & 1,469 & 1,313 & 1,208 & 1,001 \\
\quad NDX +5m & & 1,457 & 1,305 & 1,200 & 993 \\
\quad NDX +15m & & 1,435 & 1,283 & 1,178 & 972 \\
\quad NDX +60m & & 1,350 & 1,200 & 1,099 & 910 \\
\bottomrule
\end{tabular}
\begin{tablenotes}
\footnotesize
\item \textit{Notes.} H1 event/baseline magnitude ratio at four
cluster-gap thresholds. The 15-minute column corresponds to the
specification used in the main body. Tighter thresholds (5--10 min)
retain more independent clusters and yield slightly higher ratios;
a wider threshold (30 min) compresses the sample but the central
finding survives in every cell.
\end{tablenotes}
\end{threeparttable}
\end{table}


\subsection{Subsample analysis by LLM sentiment label}
\label{app:subsample}

The bull-versus-bear asymmetry test (H12) reported as null in
Section~\ref{sec:limitations} aggregates over all events.
Tables~\ref{tab:subsample-magnitude} and~\ref{tab:subsample-direction}
report the magnitude and direction results stratified by the LLM
sentiment label (bull, bear, neutral). The magnitude effect is
broadly symmetric across the three subsamples and confirms the H12
null finding at a more granular level. The direction hit rate by
sentiment label shows one notable cell: EUR/USD events labelled
bull achieve a $57.3\%$ hit rate at the $+60$-minute window
($p = 0.0007$ on a one-sided binomial); the corresponding bear
subsample is not statistically distinguishable from chance. We
report this observation as exploratory; it does not feature in the
main results because the same conclusion would not hold under a
stricter joint-testing correction across the eighteen subsample
cells reported here.

\begin{table}[!htbp]
\centering
\caption{Magnitude by sentiment label (subsample analysis).}
\label{tab:subsample-magnitude}
\begin{threeparttable}
\small
\begin{tabular}{lccccc}
\toprule
Asset & Window & Sentiment & $N$ & Mean $|r|$ (\%) & Median $|r|$ (\%) \\
\midrule
EUR/USD & +5m & bull & 557 & 0.0333 & 0.0207 \\
EUR/USD & +5m & bear & 464 & 0.0330 & 0.0207 \\
EUR/USD & +5m & neutral & 384 & 0.0306 & 0.0188 \\
EUR/USD & +15m & bull & 557 & 0.0547 & 0.0346 \\
EUR/USD & +15m & bear & 464 & 0.0527 & 0.0360 \\
EUR/USD & +15m & neutral & 384 & 0.0518 & 0.0311 \\
EUR/USD & +60m & bull & 557 & 0.1006 & 0.0677 \\
EUR/USD & +60m & bear & 464 & 0.1025 & 0.0714 \\
EUR/USD & +60m & neutral & 384 & 0.0935 & 0.0630 \\
NDX & +5m & bull & 475 & 0.0816 & 0.0506 \\
NDX & +5m & bear & 649 & 0.0923 & 0.0554 \\
NDX & +5m & neutral & 281 & 0.0700 & 0.0458 \\
NDX & +15m & bull & 475 & 0.1409 & 0.0933 \\
NDX & +15m & bear & 649 & 0.1471 & 0.0931 \\
NDX & +15m & neutral & 281 & 0.1430 & 0.0866 \\
NDX & +60m & bull & 475 & 0.2619 & 0.1606 \\
NDX & +60m & bear & 649 & 0.2806 & 0.1745 \\
NDX & +60m & neutral & 281 & 0.2539 & 0.1751 \\
\bottomrule
\end{tabular}
\begin{tablenotes}
\footnotesize
\item \textit{Notes.} Event-window absolute return $|r|$ stratified
by LLM sentiment label. The magnitude effect is broadly symmetric
across bull, bear, and neutral subsamples; the small differences
do not motivate a separate H12-style asymmetry claim.
\end{tablenotes}
\end{threeparttable}
\end{table}

\begin{table}[!htbp]
\centering
\caption{Direction hit rate by LLM sentiment label.}
\label{tab:subsample-direction}
\begin{threeparttable}
\small
\begin{tabular}{lcccccc}
\toprule
Asset & Window & Sentiment & $N$ & Hit rate & CI95 & $p$-binom \\
\midrule
EUR/USD & +5m & bull & 523 & 50.1\% & [0.458, 0.544] & 0.500 \\
EUR/USD & +5m & bear & 475 & 54.9\%$^{**}$ & [0.505, 0.594] & 0.017 \\
EUR/USD & +15m & bull & 521 & 52.4\% & [0.481, 0.567] & 0.147 \\
EUR/USD & +15m & bear & 470 & 52.6\% & [0.480, 0.570] & 0.144 \\
EUR/USD & +60m & bull & 494 & 57.3\%$^{***}$ & [0.529, 0.616] & 0.0007 \\
EUR/USD & +60m & bear & 450 & 50.4\% & [0.458, 0.550] & 0.444 \\
NDX & +5m & bull & 441 & 54.0\%$^{*}$ & [0.493, 0.586] & 0.053 \\
NDX & +5m & bear & 560 & 51.8\% & [0.476, 0.559] & 0.211 \\
NDX & +15m & bull & 434 & 53.5\%$^{*}$ & [0.488, 0.581] & 0.082 \\
NDX & +15m & bear & 551 & 49.7\% & [0.456, 0.539] & 0.568 \\
NDX & +60m & bull & 403 & 53.8\%$^{*}$ & [0.490, 0.587] & 0.067 \\
NDX & +60m & bear & 515 & 46.2\% & [0.420, 0.505] & 0.961 \\
\bottomrule
\end{tabular}
\begin{tablenotes}
\footnotesize
\item \textit{Notes.} Direction hit rate separately for events
labelled bull and bear by the LLM (neutral events excluded as
non-directional). Bull-labelled EUR/USD events at $+60$ minutes
reach $57.3\%$ hit rate ($p = 0.0007$). $^{*}p<0.10$, $^{**}p<0.05$,
$^{***}p<0.01$ on raw $p$-values.
\end{tablenotes}
\end{threeparttable}
\end{table}


\subsection{LLM prompt full text}
\label{app:llm-prompt}

The production LLM call uses the system prompt below. The
benchmark-adapted prompt for Financial PhraseBank
(Section~\ref{ssec:sentiment-validation}) is identical to this
prompt up to the schema (positive/negative/neutral in place of the
USD/NDX bull/bear/neutral fields) and the omission of the asset-
specific instruction.

\begin{quote}\small
\texttt{You are a financial sentiment analyst. Classify how a news
event likely affects the US Dollar (USD) and the Nasdaq-100 index
(NDX) in the very short term (minutes to hours after publication).
Output STRICTLY a single JSON object with this exact schema:
sentiment\_usd $\in$ \{bull, bear, neutral\}, sentiment\_ndx in
the same support, signed directional\_strength\_usd and
directional\_strength\_ndx in $[-1, +1]$, expected\_magnitude in
\{low, med, high\}, surprise\_level in \{expected, surprise,
shock\}, confidence in $[0, 1]$, and a short rationale.}

\smallskip

\texttt{Rules: base judgment strictly on the supplied text; do not
invent context; if ambiguous or off-topic, return neutral with
low confidence (\textless 0.4); USD and NDX are different assets
and their sentiments can differ (e.g., risk-off events: bull USD,
bear NDX); directional\_strength values respect sentiment direction;
return ONLY the JSON object, no preamble.}
\end{quote}

\noindent
Eight few-shot examples covering geopolitical escalation, dovish and
hawkish central-bank communication, neutral or low-relevance events,
corporate earnings, and commodity/geopolitical shocks anchor the
format. The API call uses temperature $0.1$.
